\scanpage{7}

% archive sheet number: 6

\diagram{\begin{tikzcd}[ampersand replacement=\&]
B\otimes C\& C\arrow[l,"q_C"']\\
C\arrow[u,"{\mathrm{pr}_2}"]\&D\arrow[l,"h"']\arrow[u,"h"']
\end{tikzcd}}

\deleted{i.e. \unsure{$B$}}

est-il vrai qu'on peut factoriser $h$ \deleted{\unsure{p}} à l'aide de $h':D\to A$ (une augmentation) \deleted{que} par la formule

\[
h(x)=h'(x).1_C\qquad ??
\]

La commutativité du diagramme ci-dessus signifie simplement que $h$ prend ses valeurs dans le sous-ensemble $C'$ de $C$ \deleted{formé des $x\in$} défini par

\[
x\in C'\quad\Longleftrightarrow\quad\deleted{\unsure{q_C x|p_2}\;1\otimes x=}\ q_C x=1_B\otimes x
\]

i.e. les \uline{éléments} \inserted{\deleted{\unsure{quasi}} \uline{relativement}} \uline{primitifs} de $C$ (\unsure{v.z|viz|p.r.|rel.} : $B$) C'est là une sous-$A$-algèbre de $C$, et on voit qu'en tout état de cause (\unsure{même} si $C$ n'est pas fibré principal)

\uline{le quotient de} $C$ \uline{par le groupe d'opérateurs} $B$ \uline{dans la catégorie} $\unsure{\mathcal{C}|\mathscr{C}}$ \uline{existe}, \uline{et est donné par la sous-algèbre} $C'$ de $C$ \uline{formé des éléments} \deleted{\unsure{quasi}} \inserted{\uline{relativement}} \uline{primitifs} de $C$ (\unsure{opér.|p.r.|par.} : $B$).
